\begin{ExerciseBlock}
\ExerciseHeading{03}{Rollende Vollkugel mit Rollreibung}{rigid_bodies_sphere_rolling_up_an_incline}
Eine Vollkugel mit Radius $R$ und Masse $m$ rollt mit der Geschwindigkeit $v_0 = 5,73\,$m/s auf den Fußpunkt einer schiefen Ebene ($\alpha = 5^\circ$) zu und rollt dann bis zur Höhe $h$ hinauf. Modellieren Sie die Rollreibung als Kraft $F_\mathrm{R}=\mu F_\mathrm{N}$ mit $\mu=0{,}022$.

\begin{enumerate}
	\item Bis zu welcher Höhe $h$ rollt die Kugel?

	\item Welcher Bruchteil der ursprünglichen mechanischen Gesamtenergie wird beim Hinaufrollen in Wärme $\Delta Q$ umgesetzt?

\end{enumerate}

\begin{center}
	\CachedTikzFigure{content/resources/chapters/rigid_bodies/tikz/sphere_rolling_up_an_incline.tex}
\end{center}
\end{ExerciseBlock}
\UseExerciseSolution{rigid_bodies_sphere_rolling_up_an_incline}
